\section{Broader impact and governance}
False scientific-authority promotion can amplify incorrect or misattributed claims, particularly when a system presents verification artifacts that appear independent but are not. A fail-closed layer may reduce this risk, but excessive blocking can delay legitimate scientific communication or concentrate power in whoever controls evidence and evaluator registries. Work on academic integrity and agentic abstention reinforces that refusal can be the correct behavior under genuine insufficiency, while also showing that productive abstention remains difficult and should not be conflated with generic non-response \citep{sciintegrity2026,agentabstain2026}. We therefore report abstention cost, clean coverage, custody identity, and artifact transparency rather than optimizing only a security metric.

Protected evaluation creates another trade-off: some labels must remain hidden during scoring, yet publication requires enough evidence for audit. The review supplement therefore exposes safe aggregates, a post-evaluation public slice, and a separate-implementation check while protected holdout labels and raw traces remain in restricted custody. This supports local implementation separation, not external result-verifier custody.
